\documentclass[11pt]{article}
\usepackage[margin=2.5cm]{geometry}
\usepackage{amsmath,amssymb}
\usepackage{booktabs}
\usepackage[hidelinks]{hyperref}

\newcommand{\het}{\nu}
\newcommand{\pd}{\partial}
\newcommand{\Op}{\operatorname{Op}^w}
\newcommand{\etal}{\eta_{\ell}}

\title{The Terrain-Following One-Way Equation as a Pseudodifferential Problem}
\author{Note accompanying the TAP 2026 manuscript (\texttt{main.tex})}
\date{}

\begin{document}
\maketitle

\begin{abstract}
The Liouville-transformed equation $\varphi_{\xi\xi}+B\varphi=0$ (equation
(11) of the manuscript) is read in the language of pseudodifferential
operators: $\xi$ is the evolution variable and $B$ is an operator in $\eta$
with $\xi$-dependent coefficients.
\begin{itemize}
\item We compute the exact Weyl symbol of $B$. Its principal part is the local
  dispersion relation, and it has no first-order part. The only corrections are
  two zeroth-order terms of relative size $(\lambda/\ell)^2$: the potential $V$
  of the manuscript, and a term $r_\eta^2/(4r^4)$ that the operator form keeps
  implicitly.
\item The one-way factorization (13) is the symbol expansion of the operator
  Riccati equation. We give its next term, which reduces to the classical
  second-order WKB correction.
\item The Bremmer coupling series makes the neglected backward wave explicit.
  Its local part reproduces the second-order term. Genuine backscatter is
  $O(\varepsilon^\infty)$, exponentially small for smoothed terrain. The flux
  pairs are the interface transmission coefficients without the staircase
  reflection loss.
\item The flux-conserving step, the rational propagator, the per-class
  reference wavenumber, the dispersion analysis, the vertical stretching and the
  nonuniform range step are then statements about symbols.
\item The same calculus explains why the shift map has an error floor that does
  not vanish under refinement, and why no local top boundary absorbs grazing
  waves.
\end{itemize}
\end{abstract}

\section{Setting and conventions}

\paragraph{The equation.} In orthogonal coordinates $(\xi,\eta)$ with scale factors
$h_\xi,h_\eta$ and anisotropy $r=h_\eta/h_\xi$, the Helmholtz equation reads
(manuscript (7))
\begin{equation}\label{eq:helm}
\pd_\xi(r\pd_\xi\psi)+\pd_\eta(r^{-1}\pd_\eta\psi)+h_\xi h_\eta k^2\psi=0 .
\end{equation}
With $\psi=r^{-1/2}\varphi$ it becomes (manuscript (11))
\begin{equation}\label{eq:liou}
\varphi_{\xi\xi}+B\varphi=0,\qquad
B\varphi=r^{-1/2}\pd_\eta\!\left(r^{-1}\pd_\eta(r^{-1/2}\varphi)\right)+(k^2h_\xi^2-V)\varphi,
\end{equation}
with $V=(r^{1/2})_{\xi\xi}/r^{1/2}$.

\paragraph{Quantization.} Let $D=-i\pd_\eta$ and let $q$ be the variable dual
to $\eta$; $q$ is the computational vertical wavenumber of Section VII of the
manuscript. For a function $a(\eta,q)$, depending on $\xi$ as a parameter, the
Weyl quantization is
\[
\Op(a)u(\eta)=\frac1{2\pi}\iint e^{i(\eta-\eta')q}\,a\!\left(\tfrac{\eta+\eta'}2,q\right)u(\eta')\,d\eta'\,dq .
\]
Real symbols give symmetric operators. Composition is given by the Moyal product
\begin{equation}\label{eq:moyal}
a\#b=ab+\frac i2\{a,b\}-\frac18\left(a_{qq}b_{\eta\eta}-2a_{q\eta}b_{q\eta}+a_{\eta\eta}b_{qq}\right)+\dots,\qquad
\{a,b\}=a_qb_\eta-a_\eta b_q .
\end{equation}
Odd-order terms are antisymmetric in $(a,b)$. Hence $a\#a=a^2+O(\pd^4)$, and the
symbol of a commutator is $i\{a,b\}$ to leading order.

\paragraph{The small parameter.} Let $\ell$ be the shortest length on which the
metric varies: the smoothing length of the terrain, or the transition height
$L$. The expansion parameter is $\varepsilon=1/(k\ell)$. In rescaled variables
$\eta/\ell$ and $q/k$, every derivative of a coefficient brings one factor
$\varepsilon$; this is the semiclassical calculus with $\hbar=\varepsilon$
\cite{zworski,dimassi}. Symbols are graded by powers of $\varepsilon$. We write
$a=a_0+\varepsilon a_1+\varepsilon^2a_2+\dots$ but keep the physical
variables in the formulas.

\section{The symbol of $B$}

\paragraph{Exact Weyl symbol.} For any smooth $g$ we have
$\Op(g\,q^2)=\tfrac14(D^2g+2DgD+gD^2)$. Moreover
$D^2g+gD^2=2DgD-g''$, so
\begin{equation}\label{eq:DgD}
DgD=\Op\!\left(g\,q^2+\tfrac14 g''\right).
\end{equation}
With $f=r^{-1/2}$ and $g=r^{-1}$, moving $f$ through $D$ gives
$fDgDf=D(f^2g)D-(ff'g)'+g f'^2$. The derivative part of $B$ is
$-fDgDf$. Using \eqref{eq:DgD} for $f^2g=r^{-2}$, the second derivatives of $r$
cancel and only a square of $r_\eta$ is left:
\begin{equation}\label{eq:symbol}
\boxed{\;b(\xi,\eta,q)=k^2h_\xi^2-\frac{q^2}{r^2}-V-\frac{r_\eta^2}{4r^4}\;},
\qquad B=\Op(b)\ \text{exactly}.
\end{equation}
The identity $r^{-1/2}(r^{-1}(r^{-1/2}u)')'=(r^{-2}u')'+\tfrac14(r^{-2})''u-\tfrac14r'^2r^{-4}u$
behind \eqref{eq:symbol} was checked numerically to working precision.

\paragraph{Principal symbol: the dispersion relation.} The principal symbol is
\[
b_0=k^2h_\xi^2-\frac{q^2}{r^2}.
\]
Let $\kappa$ be dual to $\xi$. The characteristic set of
\eqref{eq:liou} is $\kappa^2=b_0$, that is,
$(\kappa/h_\xi)^2+(q/h_\eta)^2=k^2$: the local plane waves of the manuscript.
With $q=kh_\eta\sin\theta$, where $\theta$ is the angle to the level curve,
\begin{equation}\label{eq:angle}
b_0=k^2h_\xi^2\cos^2\theta,\qquad \sqrt{b_0}=kh_\xi\cos\theta .
\end{equation}

\paragraph{No first-order part.} The symbol \eqref{eq:symbol} is real, with
nothing of order $\varepsilon^1$. This is the symbol form of ``$B$ is symmetric
in $L^2(d\eta)$''. It matters: a first-order (subprincipal) term would be an
imaginary part of order $\varepsilon q$, which transports amplitude across
$\eta$ and makes $B$ non-normal. The original equation \eqref{eq:helm}, divided
by $r$, has first-order terms in both variables: $(r_\xi/r)\pd_\xi$ and
$(\pd_\eta r^{-1})/r^{-1}\cdot\pd_\eta$. The Liouville substitution is the gauge
that removes them. The price is the two zeroth-order terms of order
$\varepsilon^2$:
\begin{itemize}
\item $V=(r^{1/2})_{\xi\xi}/r^{1/2}$ from the $\xi$ direction. Since $\xi$ is
  the evolution variable and is not quantized, this term has to be written
  explicitly, and the scheme computes it by a second difference along the
  columns.
\item $W=r_\eta^2/(4r^4)$ from the $\eta$ direction. It is not written in the
  manuscript and does not need to be: the scheme discretizes the operator form
  $r\pd_\eta r^{-1}\pd_\eta$, which contains $W$ exactly.
\end{itemize}
Both are $O((\lambda/\ell)^2)$ relative to $k^2h_\xi^2$, which is the size
stated after (11) for $V$.

\paragraph{The vertical stretching.} A stretching $\eta=Z(\eta')$ with
$s=dZ/d\eta'$ multiplies $h_\eta$ by $s$. In \eqref{eq:symbol} this replaces $r$
by $rs$, so $q'/(h_\eta s)$ is the physical vertical wavenumber. This is the
symbol form of ``a stretching is a metric'' (\texttt{nonuniform\_grids.md},
Section 1).

\section{One-way factorization as a symbol expansion}

\paragraph{The Riccati equation.} Seek $\varphi_\xi=i\Lambda\varphi$ with
$\Lambda=\Op(\lambda)$. Then \eqref{eq:liou} requires
\begin{equation}\label{eq:riccati}
\Lambda^2-i\Lambda_\xi=B
\quad\Longleftrightarrow\quad
\lambda\#\lambda-i\pd_\xi\lambda=b .
\end{equation}
We expand $\lambda=\lambda_0+\lambda_1+\lambda_2+\dots$, with $\lambda_j$ of
order $\varepsilon^j$ relative to $\lambda_0$, and $\pd_\xi$ of a symbol counts
as one order.

\paragraph{Order 0.} $\lambda_0^2=b_0$. The outgoing root is
$\lambda_0=\sqrt{b_0}$, with $\operatorname{Im}\sqrt{b_0}\ge0$; for evanescent
waves ($b_0<0$) it decays in $\xi$.

\paragraph{Order 1.} The first Moyal term of $\lambda\#\lambda$ vanishes, so
$2\lambda_0\lambda_1-i\pd_\xi\lambda_0=0$ and
\begin{equation}\label{eq:lambda1}
\lambda_1=\frac i2\frac{\pd_\xi\lambda_0}{\lambda_0}=\frac i4\frac{\pd_\xi b}{b}.
\end{equation}
This is the term $\tfrac i4B^{-1}B_\xi$ of (13) in the manuscript, and the
manuscript's ``treated as commuting'' is exact at this order. Indeed,
$\tfrac i4B^{-1}B_\xi$ has the symbol
$\tfrac i4[b_\xi/b+\tfrac i2\{b^{-1},b_\xi\}+\dots]$, whose correction is of
order $\varepsilon^2$. The operator ordering only matters from the next order
on.

\paragraph{Order 2.} Collecting the terms of order $\varepsilon^2$ in
\eqref{eq:riccati}, including the second Moyal term of $\lambda_0\#\lambda_0$,
gives
\begin{equation}\label{eq:lambda2}
\lambda_2=\frac1{2\sqrt b}\left[\frac{5\,b_\xi^2}{16\,b^2}-\frac{b_{\xi\xi}}{4b}
+\frac14\left(\pd_q^2\lambda_0\,\pd_\eta^2\lambda_0-(\pd_q\pd_\eta\lambda_0)^2\right)\right].
\end{equation}
\begin{itemize}
\item If $b$ does not depend on $\eta$, the last bracket vanishes. What remains
  is the classical second-order WKB correction
  $\tfrac5{32}Q'^2Q^{-5/2}-\tfrac18Q''Q^{-3/2}$ for $y''+Qy=0$: a check of the
  expansion.
\item The bracket is the phase-space curvature of the one-way symbol. Together
  with the ordering ambiguity of $\tfrac i4B^{-1}B_\xi$ (a term $\propto
  \{b^{-1},b_\xi\}$), it is what the manuscript calls ``second-order terms and
  commutators $[B,B_\xi]$'', whose symbol is $i\{b,b_\xi\}$.
\item All of these are $O(\varepsilon^2)$. Over a transition region of length
  $\ell_\xi$ they accumulate a phase of order $k\ell_\xi\varepsilon^2$.
\end{itemize}

\paragraph{Limits of the expansion.}
\begin{itemize}
\item \emph{Smoothness.} The expansion needs $\varepsilon\ll1$. The terrain
  smoothing and the transition height are what keep the metric smooth on the
  wavelength scale.
\item \emph{Turning points.} The expansion fails where $b_0=0$, i.e.\
  $\theta\to90^\circ$ to the marching direction, because $\sqrt{b_0}$ is not
  smooth there. Backscatter is not captured; this is the general limit of
  one-way methods \cite{fishman,dehoop}. Section~\ref{sec:bremmer} quantifies it. Within the aperture, $b_0\ge
  k^2h_\xi^2\cos^2\theta_{\max}>0$ keeps the field away from it.
\end{itemize}

\section{Flux and the factors $B^{\pm1/4}$}

\paragraph{Flux.} The flux $\mathcal F=\operatorname{Re}\langle\varphi,\sqrt
B\varphi\rangle$ is the quadratic form with symbol $\sqrt{b_0}$, by (14) of the
manuscript. By \eqref{eq:angle} it is the squared amplitude times
$kh_\xi\cos\theta$, which is the WKB amplitude law of the manuscript.

\paragraph{Conjugation.} The variable $\chi=B^{1/4}\varphi$ obeys
\[
\chi_\xi=B^{1/4}\left(i\sqrt B-\tfrac14B^{-1}B_\xi\right)\varphi+\left(B^{1/4}\right)_\xi\varphi
=i\sqrt B\chi+O(\varepsilon^2),
\]
because $(B^{1/4})_\xi$ has the symbol $\tfrac14b^{-3/4}b_\xi$ up to
$O(\varepsilon^2)$. The conjugation removes $\lambda_1$ exactly. So the
transport term is a gauge of order $\varepsilon$, and the flux-conserving step
is a gauge-corrected exponential.

\paragraph{The flux pairs.} In (15) of the manuscript the middle factor
$e^{i\Delta\xi\sqrt{B_{n+1/2}}}$ commutes exactly with $B_{n+1/2}^{\pm1/4}$: all
three are functions of one operator, and for symbols $\{F(b),G(b)\}=F'G'\{b,b\}=0$.
The two brackets have the symbols
\[
\left(\frac{b_{n+1/2}}{b_{n+1}}\right)^{1/4}=1-\frac{\Delta\xi}8\frac{\pd_\xi b}{b}+O(\Delta\xi^2),\qquad
\left(\frac{b_n}{b_{n+1/2}}\right)^{1/4}=1-\frac{\Delta\xi}8\frac{\pd_\xi b}{b}+O(\Delta\xi^2).
\]
\begin{itemize}
\item They are the identity wherever the two columns coincide, in particular in
  the exterior.
\item Their composition with the exponential is exact up to Poisson brackets
  between neighbouring columns,
  $\{b_n,b_{n+1/2}\}=\tfrac{\Delta\xi}2\{b,\pd_\xi b\}+\dots$. These are the
  neglected commutators of the previous section, now per step.
\end{itemize}

\paragraph{The rotated Padé factors.} Milinazzo's rotation of the factors
$(1+L)^{\pm1/4}$ moves their poles off the evanescent part of the symbol range
$b<0$. A pair of factors at slightly different columns, whose symbols almost
cancel, then no longer has a pole at almost the same real point, which is where
the spurious amplification arose.

\section{The Bremmer coupling series}\label{sec:bremmer}

The one-way equation drops the backward wave. The Bremmer series
\cite{bremmer,atkinson,dehoop} keeps both waves and treats their coupling as a
perturbation. It shows what the one-way step neglects, and how large that is.

\paragraph{Splitting.} Let $B>0$ on the propagating part of the spectrum, with
$\Lambda=B^{1/2}$ and $Q=B^{1/4}$, so that $Q^2=\Lambda$ and $Q^{-1}BQ^{-1}=\Lambda$
exactly. Define the flux-normalized amplitudes $u_\pm$ by
\begin{equation}\label{eq:split}
\varphi=Q^{-1}(u_++u_-),\qquad \varphi_\xi=iQ(u_+-u_-),\qquad
u_\pm=\tfrac12\left(Q\varphi\mp iQ^{-1}\varphi_\xi\right).
\end{equation}
Substituting into \eqref{eq:liou} gives, with no approximation,
\begin{equation}\label{eq:coupled}
\pd_\xi\begin{pmatrix}u_+\\u_-\end{pmatrix}
=\begin{pmatrix}i\Lambda+T & \mathcal R\\ \mathcal R & -i\Lambda+T\end{pmatrix}
\begin{pmatrix}u_+\\u_-\end{pmatrix},\qquad
T=\tfrac12\left[Q_\xi,Q^{-1}\right],\quad
\mathcal R=\tfrac12\left(Q_\xi Q^{-1}+Q^{-1}Q_\xi\right).
\end{equation}
\begin{itemize}
\item $Q_\xi$ and $Q^{-1}$ are symmetric, so $T$ is antisymmetric and $\mathcal R$
  is symmetric. Hence
  $\pd_\xi\left(\|u_+\|^2-\|u_-\|^2\right)=0$: the coupled system conserves the
  flux exactly, and $\|u_\pm\|^2$ are the forward and backward fluxes.
\item The symbols follow from the Moyal product \eqref{eq:moyal}, whose
  first-order term cancels in the anticommutator and survives in the
  commutator:
  \begin{equation}\label{eq:RT}
  \rho:=\operatorname{symbol}(\mathcal R)=\frac{\pd_\xi b}{4b}\bigl(1+O(\varepsilon^2)\bigr),\qquad
  \operatorname{symbol}(T)=\frac i{32}\frac{\{b,\pd_\xi b\}}{b^2}+O(\varepsilon^3).
  \end{equation}
  Relative to $\lambda_0=\sqrt b$, the coupling $\rho$ is of order
  $\varepsilon$ and $T$ of order $\varepsilon^2$. $T$ is the commutator
  $[B,B_\xi]$ that the manuscript neglects, and it is a pure phase.
\item In the exterior, and wherever the columns do not change, $\rho=0$ and
  $T=0$, and the two waves decouple exactly.
\end{itemize}

\paragraph{The series.} Let $G_\pm(\xi,\xi')$ be the propagators of
$\pd_\xi\mp i\Lambda-T$. For flux-normalized variables they are unitary. With
the incident field given at $\xi=0$ and no backward wave entering from the far
end $\xi=X$,
\[
u_+(\xi)=G_+(\xi,0)u_+(0)+\int_0^\xi G_+(\xi,\xi')\mathcal R u_-(\xi')\,d\xi',\qquad
u_-(\xi)=-\int_\xi^X G_-(\xi,\xi')\mathcal R u_+(\xi')\,d\xi'.
\]
Iterating, starting from $u_-=0$, gives the Bremmer series. Its terms alternate
between forward and backward propagation, with one scattering by $\mathcal R$ at
each turn.
\begin{itemize}
\item The zeroth term $u_+=G_+u_+(0)$ is the one-way equation in the variable
  $\chi=B^{1/4}\varphi$ of the previous section. When $u_-=0$,
  \eqref{eq:split} gives $u_+=Q\varphi=\chi$ exactly. The flux-conserving step
  of the manuscript is the zeroth Bremmer term, with $T$ dropped.
\item Because $G_\pm$ are unitary, the $n$-th term is bounded by
  $\left(\int_0^X\|\mathcal R\|\,d\xi\right)^n$. The series therefore converges
  geometrically when $\int\|\mathcal R\|d\xi<1$. By \eqref{eq:RT} this
  integral is $\tfrac12$ of the total variation of $\log(h_\xi\cos\theta)$
  along the path, which is small for the smoothed terrains of the manuscript
  as long as $\cos\theta$ stays away from $0$.
\end{itemize}

\paragraph{Two parts of the first reflection.} Take a forward wave
$u_+\approx e^{i\Phi}$, with $\Phi_\xi=\lambda_0$, and consider one column
first. The first backward term is
\begin{equation}\label{eq:back}
u_-(\xi)=-e^{-i\Phi(\xi)}\int_\xi^X\rho(\xi')\,e^{2i\Phi(\xi')}\,d\xi' .
\end{equation}
Integration by parts splits it into a local part and a remainder:
\[
u_-(\xi)=\underbrace{\frac{\rho}{2i\lambda_0}\,e^{i\Phi(\xi)}}_{\text{slaved, local}}
\;+\;\underbrace{e^{-i\Phi(\xi)}\int_\xi^X\left(\frac{\rho}{2i\lambda_0}\right)_{\xi'}e^{2i\Phi(\xi')}\,d\xi'}_{\text{genuine backscatter}} .
\]
\begin{itemize}
\item \emph{The slaved part} travels forward with the wave: its phase is
  $e^{+i\Phi}$, not $e^{-i\Phi}$. It is a local dressing of size
  $\varepsilon$ that vanishes where the metric stops changing, and it is not a
  reflected wave. Fed back into the $u_+$ equation through $\mathcal R$, it
  changes the forward rate by $\rho\cdot\rho/(2i\lambda_0)$, so the effective
  one-way symbol becomes $\lambda_0-\rho^2/(2\lambda_0)$. Undoing the
  transformation \eqref{eq:split} adds the gauge term $-\pd_\xi(\rho/2\lambda_0)$.
  With $\rho=b_\xi/(4b)$ this gives
  \[
  -\frac{\rho^2}{2\lambda_0}-\pd_\xi\frac{\rho}{2\lambda_0}
  =\frac{5\,b_\xi^2}{32\,b^{5/2}}-\frac{b_{\xi\xi}}{8\,b^{3/2}},
  \]
  which is exactly $\lambda_2$ of \eqref{eq:lambda2} without $\eta$ dependence.
  So the second-order terms of the one-way expansion are the near-diagonal
  part of the double reflection. They are local and of order $\varepsilon^2$,
  and they can be added to the one-way operator without any backward
  propagation.
\item \emph{The genuine backscatter} is an oscillatory integral with phase
  $2\Phi$ and $\Phi_\xi=\lambda_0\ge kh_\xi\cos\theta_{\max}>0$. It has no
  stationary point. Each further integration by parts gains a factor
  $\varepsilon$, so it is $O(\varepsilon^\infty)$ for a smooth metric. If the
  metric is analytic in a strip of width $\delta$ around the real $\xi$ axis,
  it is exponentially small, about $e^{-2kh_\xi\cos\theta\,\delta}$.
\end{itemize}

\paragraph{Check.} For $y''+b(x)y=0$ with $b=k^2(1+a\tanh(x/\ell))$ and
$a=0.3$, the exact reflection coefficient and the first Bremmer term
$\left|\int\rho\,e^{2i\Phi}dx\right|$ are:
\begin{center}
\begin{tabular}{@{}lllll@{}}
\toprule
$k\ell$ & $0.5$ & $1$ & $2$ & $3$ \\
\midrule
$|R|$ exact & $1.07\cdot10^{-1}$ & $4.45\cdot10^{-2}$ & $4.44\cdot10^{-3}$ & $3.55\cdot10^{-4}$ \\
first Bremmer term & $1.07\cdot10^{-1}$ & $4.43\cdot10^{-2}$ & $4.41\cdot10^{-3}$ & $3.53\cdot10^{-4}$ \\
\bottomrule
\end{tabular}
\end{center}
The two agree to within $0.6\%$. The decay is exponential in $k\ell$, at about
the rate $e^{-\pi k_{\min}\ell}$ set by the poles of $\tanh$ at distance
$\pi\ell/2$ from the real axis.

\paragraph{Multidimensional picture.} In phase space, $\mathcal R$ couples the
forward sheet $\kappa=+\lambda_0(\eta,q)$ to the backward sheet
$\kappa=-\lambda_0(\eta,q)$. Within the aperture the two sheets are separated
by $2\lambda_0\ge2kh_-\cos\theta_{\max}$. An operator with a smooth symbol
cannot transfer energy between disjoint characteristic sets beyond
$O(\varepsilon^\infty)$, which is the microlocal form of the non-stationary
phase above. The sheets meet only where $b_0=0$. There $\rho=b_\xi/(4b)$ and
$1/\lambda_0$ blow up, and the series stops being a perturbation series. This
is the turning-point limit of the one-way method again, now as a statement
about coupling.

\paragraph{What the one-way step neglects.} The Bremmer series ranks the
neglected terms:
\begin{enumerate}
\item the phase $T$ and the slaved coupling, $O(\varepsilon^2)$, both local.
  These are the second-order terms and commutators named in the manuscript,
  and they could be included in a one-way step;
\item genuine backscatter: $O(\varepsilon^\infty)$, exponentially small for
  analytic (for example Gaussian-smoothed) terrain.
\end{enumerate}
The smoothing of the terrain is therefore not just a numerical convenience. A
kink in $T(x)$ makes the metric non-analytic. The backscatter then decays
only algebraically in $k$: it is edge diffraction, which no one-way method
represents. This is the generalized Bremmer series of de~Hoop
\cite{dehoop}. Choosing $\Lambda$ to include higher terms of the Riccati
expansion pushes $\mathcal R$ to higher order, and in the limit the forward
and backward waves are decoupled to all orders.

\paragraph{The flux pairs as interface transmission.} The column-by-column
step replaces the continuous metric by a staircase in $\xi$. At an interface
between columns with roots $\lambda_a$ and $\lambda_b=\lambda_a(1+\delta)$,
continuity of $\varphi$ and $\varphi_\xi$ gives the amplitude transmission
$t=2\lambda_a/(\lambda_a+\lambda_b)$ and reflection
$\delta/(2+\delta)$. In flux-normalized form,
\[
t=\left(\frac{\lambda_a}{\lambda_b}\right)^{1/2}\tau=\left(\frac{b_a}{b_b}\right)^{1/4}\tau,\qquad
\tau=\frac{2\sqrt{\lambda_a\lambda_b}}{\lambda_a+\lambda_b}=1-\frac{\delta^2}8+O(\delta^3),
\]
with $\tau^2+|\delta/(2+\delta)|^2=1$.
\begin{itemize}
\item The flux pairs of (15) are exactly the factor $(b_a/b_b)^{1/4}$, so the
  step is the staircase transmission with $\tau$ set to $1$.
\item This is the better choice. Here $\delta\approx2\rho\,\Delta\xi$, so each
  interface reflects a flux of about $\rho^2\Delta\xi^2$. Over a length
  $\ell_\xi$ the staircase would remove a forward flux of order
  $\rho^2\ell_\xi\Delta\xi$ through backscatter. That loss is an artefact of
  the staircase: it vanishes as $\Delta\xi\to0$, and the reflections of the
  continuous metric cancel by phase as shown above.
\item Taking $\tau=1$ removes the artefact and conserves the forward flux
  exactly. This is the discrete counterpart of the zeroth Bremmer term.
\end{itemize}

\paragraph{Caveat: evanescent part.} For $b<0$ the roots $\Lambda$ and $Q$ are
complex and $T$, $\mathcal R$ lose their (anti)symmetry. The splitting
\eqref{eq:coupled} still holds algebraically with the decaying branch, but it
no longer carries a flux interpretation. Evanescent components do not
propagate, so the estimates above concern the propagating sector.

\paragraph{Can the existing JAX propagators implement this series?}
Yes, as a computational framework, but not by simply running the existing
one-way solver forward and backward. The repository's
\texttt{RationalHelmholtzPropagator} and \texttt{SSFHelmholtzPropagator}
advance a single reduced field in range; their \texttt{compute} methods do
not expose the transverse Helmholtz operator $B$, its square root, or a
backward amplitude. In particular, reversing the range loop is not the
inverse of a lossy, terrain-masked step and would not impose the two-end
boundary conditions in the Bremmer integral equations.

The reusable parts are the JAX arrays, terrain and heterogeneity sampling,
vertical discretizations, boundary-condition machinery, and compiled range
loops. A first implementation for the existing flat-coordinate 2D problem
can form a discrete transverse operator $B_j$ at each range column using the
same vertical stencil (rational solver) or spectral multiplier (SSF solver).
On its propagating subspace, apply $B_j^{1/2}$ and $B_j^{1/4}$ with the
existing spectral representation where available, or a rational/Krylov
approximation for the finite-difference operator. Then compute the discrete
coupling from the change in $Q_j=B_j^{1/4}$, preferably as the symmetric
formula in \eqref{eq:coupled}; this avoids differentiating a noisy terrain
profile. JAX's \texttt{lax.scan}/\texttt{fori\_loop} can march $u_+$ from the
source and $u_-$ from the far-range no-incoming-wave boundary, and repeat the
two sweeps to sum a chosen number of Bremmer reflections. The far boundary
and all range-dependent terrain masks must be included in this two-point
iteration rather than inherited from the one-way absorber.

This is a moderate new solver layer, not a small option on \texttt{compute}.
The most direct prototype is SSF with a flat, lossless domain and a smooth
range-varying medium: its vertical spectrum makes the roots and flux
normalization accessible. The rational finite-difference solver can follow
once a stable matrix-function action and compatible flux form are specified.
The current top absorbing layer and lossy/complex media also need separate
care, since the unitary $G_\pm$ and conserved flux used in the convergence
bound apply only to the self-adjoint propagating problem. Validate first
against a range-dependent 1D profile with an exact reflection coefficient,
then compare the zeroth and first Bremmer orders against the existing
one-way result. Thus the existing JAX infrastructure is suitable, but a
faithful two-way implementation requires explicit discrete factorization and
a two-boundary iteration.

\section{The rational step: functional calculus}

\paragraph{Symbol of the step.} For $B=\Op(b)$ symmetric and a smooth function
$F$, the semiclassical functional calculus gives
\begin{equation}\label{eq:funcalc}
F(B)=\Op\!\left(F\circ b\right)+O(\varepsilon^2)
\end{equation}
\cite{helffer,dimassi}. The $O(\varepsilon)$ correction vanishes in the Weyl
calculus. Apply this to $F=R$, the rational step (3), in the variable
$L=B/\beta^2-1$:
\[
\operatorname{symbol}\bigl(R(L)\bigr)=R\!\left(\frac{b}{\beta^2}-1\right)+O(\varepsilon^2),
\qquad
\operatorname{symbol}\bigl(e^{i\beta\Delta\xi(\sqrt{1+L}-1)}\bigr)=e^{i\Delta\xi(\sqrt b-\beta)}+O(\varepsilon^2).
\]
So the error of the rational approximation is, to leading order, a pointwise
function on phase space,
\[
E(\xi,\eta,q)=\left|R\!\left(\tfrac{b}{\beta^2}-1\right)-e^{i\Delta\xi(\sqrt b-\beta)}\right| .
\]

\paragraph{The error budget.} What has to be small is $E$ on the support of
the field in phase space: its wavefront set, or in practice the support of its
Wigner function. By \eqref{eq:angle}, for angles up to the aperture
$\theta_{\max}$ the symbol $b$ ranges over
$[k_{\min}^2h_-^2\cos^2\theta_{\max},\,k_{\max}^2h_+^2]$, with
$h_\pm=\max/\min h_\xi$. These are the rescaled bounds with which the grid
optimizer is called (Section VII of the manuscript). The ``local
equivalence'' is the statement that $E$ depends on the metric only through this
range.

\paragraph{Per-class reference wavenumber.} On one column, $b$ takes values only
in that column's $h_\xi$ range. With $\beta_k$ chosen per class, the Padé
approximant has to cover a narrower interval of $b/\beta_k^2-1$
(\texttt{beta\_mode='column'}). The envelope factor
$e^{i(\beta_k-\beta)\Delta\xi}$ is a constant phase and does not change the
symbol.

\paragraph{Where the field lives in phase space: bicharacteristics.} The
support of the field is not a property of one column. It is transported by the
Hamiltonian flow of the one-way symbol,
\begin{equation}\label{eq:rays}
\frac{d\eta}{d\xi}=-\pd_q\sqrt{b_0},\qquad \frac{dq}{d\xi}=\pd_\eta\sqrt{b_0},
\end{equation}
and reflected at the ground by the boundary symbol (manuscript (21)).
\begin{itemize}
\item A ray reflected by a slope $\theta_s$ leaves with up to
  $\theta_{\max}+2\theta_s$ to the horizontal.
\item Once the terrain flattens, its local symbol range is that of a flat
  column, but its $q$ is still large.
\item This is why the nonuniform range step (\texttt{range\_steps='class'})
  failed when the aperture did not include the slopes. Each class's step was
  sized for the $q$ range of its own columns, while the field there carried
  the larger $q$ produced upstream. The error budget has to be evaluated on
  the transported support, i.e.\ with the slope-widened aperture
  (\texttt{nonuniform\_grids.md}, Section 7).
\end{itemize}

\section{Discrete symbols}

\paragraph{The discrete symbol of the scheme.} The discrete scheme has a
discrete symbol in $q\in(-\pi/\Delta\eta,\pi/\Delta\eta]$. For frozen
coefficients it is manuscript (18),
\[
L_h(q)=\het-\frac4{\tau^2r^2}\frac{\sigma}{1-4\alpha\sigma},\qquad \sigma=\sin^2\frac{q\Delta\eta}2,
\]
which replaces $b/\beta^2-1=\het-q^2/(\beta r)^2$ by a periodic function that
agrees with it to $O((q\Delta\eta)^6)$ for $\alpha=1/12$. The step's symbol is
$R(L_h(q))$, and the dispersion error (19) is the distance between this and the
exact symbol of the previous section.

\paragraph{Local equivalence.} Manuscript (20) says that the full discrete
symbol depends on the metric only through the products $h_\xi\Delta\xi$,
$h_\eta\Delta\eta$ and $\beta/h_\xi$. That is a statement about symbols: in
physical variables, the discrete symbol of a distorted cell equals the flat
discrete symbol on the effective grid.

\paragraph{Variable coefficients.} The weighted second difference with
half-node values $r_{j\pm1/2}$ is the discrete analogue of the operator form in
\eqref{eq:liou}. Its discrete Weyl symbol has the same structure as
\eqref{eq:symbol}: no first-order part, and a zeroth-order correction that
converges to $W$ as $\Delta\eta\to0$. This is the discrete form of the remark
that $W$ is carried implicitly.

\paragraph{The vertical step rule.} The error of $L_h$ is a function of
$q\Delta\eta=k\sin\theta\cdot h_\eta s\Delta\eta$, by \eqref{eq:angle} with
$h_\eta\to h_\eta s$. Bounding it at every level requires
$\max_\xi h_\eta s\,\Delta\eta\le\Delta z_{\rm phys}$. This is the interior
stretching rule (\texttt{nonuniform\_grids.md}, Section 6). The manuscript's
uniform rule bounds the same quantity by the cruder $h_+\max r$.

\section{The shift map in symbols}

Write the shift-map equation as
$\psi_{xx}-2s\psi_{x\zeta}+(1+s^2)\psi_{\zeta\zeta}+k^2\psi=0$, with $s=T'$. Its
symbol in $(\kappa,q)$, the duals of $(x,\zeta)$, is
$-\kappa^2+2s\kappa q-(1+s^2)q^2+k^2$, so the outgoing root is
\begin{equation}\label{eq:shift}
\kappa=sq+\sqrt{k^2-q^2}.
\end{equation}
\begin{itemize}
\item \emph{The obstruction.} The one-way symbol has a part $sq$ that is odd
  in $q$. No function of a single symmetric second-order vertical operator
  without a first-order term has a symbol with an odd part. This is the symbol
  form of ``no function of a single vertical operator is the exact propagator''
  (manuscript, Section III).
\item \emph{The gauge.} The gauge (23) multiplies by $e^{ik_0s\zeta}$, which
  translates the dual variable, $q=q'+k_0s$. Combined with the range phase it
  gives
  \[
  \kappa'=sq'+\tfrac{k_0s^2}2+\sqrt{k^2-(q'+k_0s)^2}-k_0 .
  \]
  If the square root is truncated at second order, the narrow-angle (NA)
  symbol, the terms linear in $q'$ and the constants cancel and
  $\kappa'=-q'^2/(2k_0)$ exactly: the BT transformation is exact for the
  narrow-angle PE. The higher terms of the root involve $(q'+k_0s)^4,\dots$ and
  do not cancel.
\item \emph{The floor.} A wide-angle PE uses the symbol $k_0[f(-q'^2/k_0^2)-1]$,
  an even function of $q'$, while the exact $\kappa'$ is not even beyond
  second order. The difference is the floor (25). It is a mismatch of symbols,
  independent of the grid, so no refinement removes it.
\item \emph{The orthogonal map.} The exact one-way symbol is $\sqrt{b_0}$, a
  function of the symbol of the single operator $B$, so no floor arises.
\end{itemize}

\section{The top boundary}

\paragraph{Continuous exterior.} Above the transition height, $h_\xi=h_\eta=1$
and the exterior is invariant in $x$. In the Fourier variable $\kappa$ (dual to
$x$), the field there is $e^{i\sqrt{k^2-\kappa^2}\,z}$, so the transparent
condition is
\[
\pd_z\psi=i\,\Op_x\!\left(\sqrt{k^2-\kappa^2}\right)\psi .
\]
This is a pseudodifferential operator in the evolution variable $x$, hence a
convolution in range. That is why the exact NLBC (5) carries the whole history.
Its symbol has branch points at $\kappa=\pm k$, i.e.\ at vertical wavenumber
$k_z=0$: waves that graze the top.

\paragraph{Discrete exterior.} In the scheme, the $z$-transform variable $t$
plays the role of $e^{i\kappa\Delta x}$, and the scalar exterior ratio
$R(\lambda)=u_{N-1}/u_N$ is the discrete symbol of the condition. Here
$\lambda=\beta^2(1+\sigma)$ is the spectral variable of the $\beta$-free
operator $A$ (\texttt{nonuniform\_grids.md}, Section 8). $R(\lambda)$ has the
discrete branch point where the two roots
$\mu=e^{\pm ik_z\Delta z}$ meet, again at $k_z=0$.

\paragraph{Local conditions are rational approximations.} A boundary condition
that is local in range, such as an absorbing layer or a complex-step layer of
$M$ rows, has an exterior ratio that is a rational function of $\lambda$ with
finitely many poles. A rational function cannot approximate a square root
uniformly up to its branch point:
\begin{itemize}
\item on the band $k_z\Delta z\in[\theta_{\min},\theta_{\max}]$ the best error
  decays geometrically in $M$, with a rate that depends on
  $\log(\theta_{\max}/\theta_{\min})$ \cite{newman,idk};
\item as $\theta_{\min}\to0$ the rate goes to zero.
\end{itemize}
In the continuous picture the layer's complex thickness $\Lambda$ gives a
reflection of about $e^{-2k_z\operatorname{Im}\Lambda}$, which tends to $1$ as
$k_z\to0$ for any finite layer. The complex layer of the code is the optimal-grid
construction of \cite{idk,adgk} applied to the Numerov exterior. The observed
rows (8, 12, 16 and 34 rows for $5\cdot10^{-3}$, $2\cdot10^{-4}$,
$10^{-5}$ and $10^{-6}$ over $0.5$--$30^\circ$) show this geometric rate.

\paragraph{Why the layer serves every class.} For per-class $\beta_k$, Padé
chains and range steps, the exterior equation of every substep is
$(A-\lambda)w=0$ with the same $A$. A condition expressed as a function of
$\lambda$, i.e.\ of the symbol of $A$ alone, is therefore valid for all of
them. The convolution NLBC is not: it is a function of $t$, and the relation
between $t$ and $\lambda$ changes with the class.

\section{Summary of orders}

\begin{center}
\begin{tabular}{@{}lll@{}}
\toprule
term & symbol & order \\
\midrule
dispersion relation & $b_0=k^2h_\xi^2-q^2/r^2$ & $1$ \\
first-order (transport) part of $B$ & $0$ (Liouville gauge) & --- \\
$V$, $W$ & $(r^{1/2})_{\xi\xi}/r^{1/2}$, $r_\eta^2/(4r^4)$ & $\varepsilon^2$ \\
transport term of the one-way equation & $\tfrac i4 b_\xi/b$ & $\varepsilon$ (removed by $B^{1/4}$) \\
commutators, second-order WKB & \eqref{eq:lambda2}, $i\{b,b_\xi\}$ & $\varepsilon^2$ \\
flux pairs at neighbouring columns & $\{b_n,b_{n+1/2}\}$ & $\varepsilon\,\Delta\xi$ \\
Bremmer coupling $\mathcal R$ & $b_\xi/(4b)$ & $\varepsilon$ \\
\quad slaved part (local) & $-\rho^2/(2\lambda_0)$, part of $\lambda_2$ & $\varepsilon^2$ \\
\quad genuine backscatter & $\int\rho\,e^{2i\Phi}d\xi$ & $\varepsilon^\infty$ ($e^{-2kh\cos\theta\,\delta}$) \\
staircase reflection loss (avoided by flux pairs) & $\rho^2\Delta\xi^2$ per step & $\varepsilon^2\Delta\xi$ total \\
rational step vs exponential & $R(b/\beta^2-1)-e^{i\Delta\xi(\sqrt b-\beta)}$ & set by the optimizer \\
Numerov symbol & $L_h(q)-(b/\beta^2-1)$ & $(q\Delta\eta)^6$ \\
shift map, wide angle & odd part of $sq+\sqrt{k^2-q^2}$ & $O(1)$, grid independent \\
local top condition & rational approximation of $\sqrt{\cdot}$ near $k_z=0$ & geometric in $M$, not uniform \\
\bottomrule
\end{tabular}
\end{center}

\begin{thebibliography}{10}
\bibitem{zworski} M.~Zworski, \emph{Semiclassical Analysis}. Providence, RI: AMS, 2012.
\bibitem{dimassi} M.~Dimassi and J.~Sj\"ostrand, \emph{Spectral Asymptotics in the Semi-Classical Limit}. Cambridge Univ. Press, 1999.
\bibitem{helffer} B.~Helffer and D.~Robert, ``Calcul fonctionnel par la transformation de Mellin et op\'erateurs admissibles,'' \emph{J. Funct. Anal.}, vol.~53, pp.~246--268, 1983.
\bibitem{fishman} L.~Fishman and J.~J.~McCoy, ``Derivation and application of extended parabolic wave theories,'' \emph{J. Math. Phys.}, vol.~25, pp.~285--308, 1984.
\bibitem{dehoop} M.~V.~de~Hoop, ``Generalization of the Bremmer coupling series,'' \emph{J. Math. Phys.}, vol.~37, pp.~3246--3282, 1996.
\bibitem{bremmer} H.~Bremmer, ``The W.K.B. approximation as the first term of a geometric-optical series,'' \emph{Comm. Pure Appl. Math.}, vol.~4, pp.~105--115, 1951.
\bibitem{atkinson} F.~V.~Atkinson, ``Wave propagation and the Bremmer series,'' \emph{J. Math. Anal. Appl.}, vol.~1, pp.~255--276, 1960.
\bibitem{newman} D.~J.~Newman, ``Rational approximation to $|x|$,'' \emph{Michigan Math. J.}, vol.~11, pp.~11--14, 1964.
\bibitem{idk} D.~Ingerman, V.~Druskin and L.~Knizhnerman, ``Optimal finite difference grids and rational approximations of the square root. I. Elliptic problems,'' \emph{Comm. Pure Appl. Math.}, vol.~53, pp.~1039--1066, 2000.
\bibitem{adgk} S.~Asvadurov, V.~Druskin, M.~N.~Guddati and L.~Knizhnerman, ``On optimal finite-difference approximation of PML,'' \emph{SIAM J. Numer. Anal.}, vol.~41, pp.~287--305, 2003.
\end{thebibliography}

\end{document}
